\documentclass{article}

% if you need to pass options to natbib, use, e.g.:
%     \PassOptionsToPackage{numbers, compress}{natbib}
% before loading neurips_2019

% ready for submission
% \usepackage{neurips_2019}

% to compile a preprint version, e.g., for submission to arXiv, add add the
% [preprint] option:
%     \usepackage[preprint]{neurips_2019}

% to compile a camera-ready version, add the [final] option, e.g.:
     \usepackage[nonatbib, final]{neurips_2019}
     

% to avoid loading the natbib package, add option nonatbib:
%     \usepackage[nonatbib]{neurips_2019}

\usepackage[utf8]{inputenc} % allow utf-8 input
\usepackage[T1]{fontenc}    % use 8-bit T1 fonts
\usepackage{hyperref}       % hyperlinks
\usepackage{url}            % simple URL typesetting
\usepackage{booktabs}       % professional-quality tables
\usepackage{amsfonts}       % blackboard math symbols
\usepackage{nicefrac}       % compact symbols for 1/2, etc.
\usepackage{microtype}      % microtypography

\usepackage{float}
\usepackage{graphicx}
\usepackage{subfig}

\usepackage{cite}


\title{TEM-t as The Universal Neural Computer}

% The \author macro works with any number of authors. There are two commands
% used to separate the names and addresses of multiple authors: \And and \AND.
%
% Using \And between authors leaves it to LaTeX to determine where to break the
% lines. Using \AND forces a line break at that point. So, if LaTeX puts 3 of 4
% authors names on the first line, and the last on the second line, try using
% \AND instead of \And before the third author name.

\author{
  \textbf{Renan Monteiro Barbosa$^{1}$} \\ 
  $^{1}$University of West Florida \\ 
  \texttt{rmb54f@students.uwd.edu, \{email1, email2,\}@affiliation.edu}
}

\begin{document}

\maketitle

\begin{abstract}
The abstract paragraph should be indented \nicefrac{1}{2}~inch (3~picas) on both the left- and right-hand margins. Use 10~point type, with a vertical spacing (leading) of 11~points.  The word \textbf{Abstract} must be centered, bold, and in point size 12. Two line spaces precede the abstract. The abstract must be limited to one paragraph.

Tempus quam pellentesque nec nam aliquam sem et tortor. 
Quam elementum pulvinar etiam non quam. 
Pretium aenean pharetra magna ac placerat vestibulum lectus mauris ultrices. 
Nunc aliquet bibendum enim facilisis. 
Lorem mollis aliquam ut porttitor leo a diam sollicitudin. 
Adipiscing elit pellentesque habitant morbi. 
Feugiat sed lectus vestibulum mattis ullamcorper velit sed ullamcorper morbi.
\end{abstract}


\section{Introduction}

Advancements in the last few decades have brought us to a point at which what was thought to only be theorized by Tolman (1948) as he argued that for Humans and other animals to make complex inferences from sparse observations and rapidly integrate new knowledge to control their behavior there should exist a systematic organization of such knowledge on what could be called a cognitive map. (note, try to make it sound like the theory of computing)

But what was missing from these early studies, was a way to address the neuronal mechanisms that led information to be stored as memory. Further development has shown that place cells are part of a wider network of spatially modulated neurons, including grid, border, and head direction cells, each with distinct roles in the representation of space and spatial memory.

Bringing to surface a mechanistic basis for memory formation (Nobel prize winning Place Cells, Grid Cells, and Memory) which later has been suggested that relational memory and spatial reasoning might be related by a common mechanism (Eichenbaum and Cohen, 2014).


\subsection{Retrieval of style files}

The style files for NeurIPS and other conference information are available on the World Wide Web at  \url{http://www.neurips.cc/}. The file \verb+neurips_2019.pdf+ contains these instructions and illustrates the various formatting requirements your NeurIPS paper must satisfy.

The only supported style file for NeurIPS 2019 is \verb+neurips_2019.sty+, rewritten for \LaTeXe{}. \textbf{Previous style files for \LaTeX{} 2.09,  Microsoft Word, and RTF are no longer supported!}


\subsection{Retrieval of style files}

The style files for NeurIPS and other conference information are available on the World Wide Web at  \url{http://www.neurips.cc/}. The file \verb+neurips_2019.pdf+ contains these instructions and illustrates the various formatting requirements your NeurIPS paper must satisfy.

The only supported style file for NeurIPS 2019 is \verb+neurips_2019.sty+, rewritten for \LaTeXe{}. \textbf{Previous style files for \LaTeX{} 2.09,  Microsoft Word, and RTF are no longer supported!}


\subsection{Retrieval of style files}

The style files for NeurIPS and other conference information are available on the World Wide Web at  \url{http://www.neurips.cc/}. The file \verb+neurips_2019.pdf+ contains these instructions and illustrates the various formatting requirements your NeurIPS paper must satisfy.

The only supported style file for NeurIPS 2019 is \verb+neurips_2019.sty+, rewritten for \LaTeXe{}. \textbf{Previous style files for \LaTeX{} 2.09,  Microsoft Word, and RTF are no longer supported!}


% ###################################################################################################################
% ###################################################################################################################
% Adding a Single big Image
% ###################################################################################################################
% ###################################################################################################################
\section{Single image}

Lorem ipsum dolor sit amet, consectetur adipiscing elit, sed eiusmod tempor incidunt ut labore et dolore magna aliqua. Ut enim ad minim veniam, quis nostrud exercitation ullamco laboris nisi ut aliquid ex ea commodi consequat. Quis aute iure reprehenderit in voluptate velit esse cillum dolore eu fugiat nulla pariatur. Excepteur sint obcaecat cupiditat non proident, sunt in culpa qui officia deserunt mollit anim id est laborum. Fig.~\ref{image1} shows the NeurIPS logo.


\begin{figure}[H]
  \centering{\includegraphics[width=40mm]{images/logo-neurips.png}}
  \caption{Example of single image}
  \label{image1}
\end{figure}


% ###################################################################################################################
% ###################################################################################################################
% Adding multiple Images
% ###################################################################################################################
% ###################################################################################################################
\section{Multiple images}

Lorem ipsum dolor sit amet, consectetur adipiscing elit, sed eiusmod tempor incidunt ut labore et dolore magna aliqua. Ut enim ad minim veniam, quis nostrud exercitation ullamco laboris nisi ut aliquid ex ea commodi consequat. Quis aute iure reprehenderit in voluptate velit esse cillum dolore eu fugiat nulla pariatur. Excepteur sint obcaecat cupiditat non proident, sunt in culpa qui officia deserunt mollit anim id est laborum. Fig.~\ref{image1} shows the NeurIPS logo as well as Fig.~\ref{image2}.


\begin{figure}[H]
\centering
\subfloat[]{\includegraphics[width=40mm]{images/logo-neurips.png}} \quad \quad 
\subfloat[]{\includegraphics[width=40mm]{images/logo-neurips.png}} \quad
\caption{Examples for sub-images} \label{image2}
\end{figure}


\begin{figure}[H]
  \centering
  \begin{minipage}[b]{40mm}
    \includegraphics[width=40mm]{images/logo-neurips.png}
    \caption{Logo image}
    \label{image3}
  \end{minipage} 
  \quad \quad \quad \quad \quad
  \begin{minipage}[b]{40mm}
    \includegraphics[width=40mm]{images/logo-neurips.png}
    \caption{Logo image}
    \label{image4}
  \end{minipage}
\end{figure}

% ###################################################################################################################
% ###################################################################################################################
% Example of Citations
% ###################################################################################################################
% ###################################################################################################################
\section{Citations}

Lorem ipsum dolor sit amet, consectetur adipiscing elit, sed eiusmod tempor incidunt ut labore et dolore magna aliqua. Ut enim ad minim veniam, quis nostrud exercitation ullamco laboris nisi ut aliquid ex ea commodi consequat. Quis aute iure reprehenderit in voluptate velit esse cillum dolore eu fugiat nulla pariatur. Excepteur sint obcaecat cupiditat non proident, sunt in culpa qui officia deserunt mollit anim id est laborum. 
For example, \cite{reference_1} or \cite{reference_2} and \cite{reference_3} anre example for citations.

% ###################################################################################################################
% ###################################################################################################################
% Tables
% ###################################################################################################################
% ###################################################################################################################
\section{Tables}

Lorem ipsum dolor sit amet, consectetur adipiscing elit, sed eiusmod tempor incidunt ut labore et dolore magna aliqua. Ut enim ad minim veniam, quis nostrud exercitation ullamco laboris nisi ut aliquid ex ea commodi consequat. Quis aute iure reprehenderit in voluptate velit esse cillum dolore eu fugiat nulla pariatur. Excepteur sint obcaecat cupiditat non proident, sunt in culpa qui officia deserunt mollit anim id est laborum. The Table~\ref{sample-table} shows an example of a table.

\begin{table}[H]
  \caption{Sample table title}
  \label{sample-table}
  \centering
  \begin{tabular}{lll}
    \toprule
    \multicolumn{2}{c}{Part}                   \\
    \cmidrule(r){1-2}
    Name     & Description     & Size ($\mu$m) \\
    \midrule
    Dendrite & Input terminal  & $\sim$100     \\
    Axon     & Output terminal & $\sim$10      \\
    Soma     & Cell body       & up to $10^6$  \\
    \bottomrule
  \end{tabular}
\end{table}


\section{Conclusions}

Lorem ipsum dolor sit amet, consectetur adipiscing elit, sed eiusmod tempor incidunt ut labore et dolore magna aliqua. Ut enim ad minim veniam, quis nostrud exercitation ullamco laboris nisi ut aliquid ex ea commodi consequat. Quis aute iure reprehenderit in voluptate velit esse cillum dolore eu fugiat nulla pariatur. Excepteur sint obcaecat cupiditat non proident, sunt in culpa qui officia deserunt mollit anim id est laborum.



\subsubsection*{Acknowledgments}

Use unnumbered third level headings for the acknowledgments. All acknowledgments go at the end of the paper. Do not include acknowledgments in the anonymized submission, only in the final paper. This example was prepared by Dennis Núñez Fernández.



\bibliography{references}{}
\bibliographystyle{plain}


\end{document}
